%% kapa: exam=1 round=32 year=2021 period=2 form=A subject=appraisal_law subject_ko=감정평가관계법규 seq=28 number=28
%% source: Q-Net 감정평가사 기출문제 (cst00309) artlSeq=5208397 file=extracted/1008/5208397/A형/2021년도 제32회 감정평가사 1차 2교시 A형.pdf
%% conversion: kordoc 4.24.51 -> tools/gen1.py (auto); review=proofread-sonnet; flags=none
\begin{question}{28}
공간정보의 구축 및 관리 등에 관한 법령상 용어에 관한 설명으로 옳지 않은 것은?
\begin{choices}
  \item 공공측량과 지적측량은 일반측량에 해당한다.
  \item 연속지적도는 측량에 활용할 수 없는 도면이다.
  \item 토지의 이동(異動)이란 토지의 표시를 새로 정하거나 변경 또는 말소하는 것을 말한다.
  \item 「지방자치법」에 따라 자치구가 아닌 구를 두는 시의 시장은 지적소관청에 해당하지 않는다.
  \item 「도시개발법」에 따른 도시개발사업이 끝나 토지의 표시를 새로 정하기 위하여 실시하는 지적측량은 지적확정측량에 해당한다.
\end{choices}
\end{question}
\answer{1}
